\begin{statement}{C.2.UN.28}{Proved}{The sum of the normalization geometric genera in one torus fibre is at most one}
Let $f:U\to\Delta$ be a proper holomorphic map from a
smooth complex threefold, with connected two-torus
fibres over $\Delta^*$. Write its full central fibre
as $\sum_i m_iS_i$, and let $Z_i\to S_i^\nu$ be
smooth resolutions of the normalizations. Then
\[
                         \sum_i p_g(Z_i)\le1.
                                                               \tag{C.PG.1}
\]
The assertion has no K\"ahler, Betti-number, reducedness,
or normal-crossing hypothesis. In particular at most
one normalization component in a given fibre has a
nonzero holomorphic two-form, and that component has
geometric genus one.

The same bound holds for any collection of multiplicity-one
components in a smooth resolved finite base change of
this family. The proof compares their compact homology
cores with one actual nearby torus, preserving their
intersection forms.
\end{statement}

\begin{proof}
Choose $N$ divisible by every $m_i$, make the actual
base change $t=s^N$, normalize, and resolve the total
space and its central support, isomorphically over
$s\ne0$ and at a generic point of every strict
component. This produces a proper map from a smooth
threefold
\[
                         g:\widetilde U\to\Delta_s
                                                               \tag{C.PG.2}
\]
whose nearby fibre is exactly $f^{-1}(s^N)$.
Every strict component above an $S_i$ has multiplicity
one. To see the map explicitly at a generic point,
write $t=u x^{m_i}$ with $u$ a unit. The normalized
root charts are
\[
       w^{m_i}=u,\qquad x=s^{N/m_i}/w.
                                                               \tag{C.PG.3}
\]
The unit-root equation is nonsingular at $s=0$;
there $s$ is a normal parameter. Subsequent
resolution is an isomorphism at this generic
point and preserves the multiplicity one.

Select one smooth strict component $E_i$ above each
$S_i$. Its map to $S_i^\nu$ is generically finite.
Resolve the actual rational comparison to $Z_i$:
\[
              E_i\xleftarrow{\rho_i}W_i
                            \xrightarrow{v_i}Z_i.
                                                               \tag{C.PG.4}
\]
The first map is a sequence of point blowups,
and the second is generically finite. Pullback
of holomorphic two-forms through $v_i$ is
injective, since its Jacobian is nonzero on
a dense open set. Point-blowup invariance of
holomorphic two-forms follows by restriction
and Hartogs extension across the blown-down
points. Hence
\[
                         p_g(E_i)\ge p_g(Z_i).
                                                               \tag{C.PG.5}
\]
It suffices to prove the bound for the selected
components $E_i$ of (C.PG.2).

Let $D_i$ be the reduced intersection of $E_i$
with all the other central components of the
resolved support. It is a compact normal-crossing
curve, or is empty. The restriction of $g$ is
a submersion near $E_i\setminus D_i$: in the
normal-crossing local equation only this
multiplicity-one component remains, so $g$
is a unit times its local defining coordinate.

We first specify the relevant homology map on
any such compact smooth surface $E=E_i$ and
boundary curve $D=D_i$. Poincar\'e--Alexander
duality, applied to the triangulable compact
curve $D$ in the oriented real four-manifold
$E$, identifies the following portion of the
relative homology sequence:
\[
 H_2(E\setminus D,\mathbb R)\longrightarrow
 H_2(E,\mathbb R)\xrightarrow{\ I_D\ }
 H_2(E,E\setminus D;\mathbb R)
       \simeq H^2(D,\mathbb R).
                                                               \tag{C.PG.6}
\]
There is one basis vector in $H^2(D,\mathbb R)$
for each irreducible curve of $D$. In these
coordinates $I_D$ is precisely the vector of
intersection numbers with those curves: under
Poincar\'e duality the map is restriction of
$H^2(E)$ to $D$. Thus its kernel is exactly
the image from the open complement in
(C.PG.6). This statement includes all
intersections between boundary curves; no
disjointness of those curves is assumed.

Every holomorphic two-form $\alpha$ on $E$
is closed, because both its $(3,0)$ derivative
and its antiholomorphic derivative vanish.
The real classes of these forms give a
positive subspace
\[
 P_E=\{[\operatorname{Re}\alpha]:
                       \alpha\in H^0(E,K_E)\}
       \subset H^2(E,\mathbb R),
 \qquad \dim_{\mathbb R}P_E=2p_g(E).
                                                               \tag{C.PG.7}
\]
Indeed, with the complex orientation,
\[
 \int_E(\operatorname{Re}\alpha)^2
       =\tfrac12\int_E\alpha\wedge\overline\alpha>0
                      \quad (\alpha\ne0).
                                                               \tag{C.PG.8}
\]
This also proves injectivity of the real-part
map on the underlying real vector space of
holomorphic forms. Every such class has zero
integral on every boundary curve, since the
pullback of a two-form to the normalization
of a curve is zero. Its Poincar\'e-dual
homology class therefore lies in
$\ker I_D$ and lifts to
$H_2(E\setminus D,\mathbb R)$ by (C.PG.6).
No Hodge decomposition or K\"ahler form has
been used.

Choose a real basis of each $P_{E_i}$,
choose lifts of its Poincar\'e-dual classes,
and represent those lifts by finite compact
smooth singular cycles in $E_i\setminus D_i$.
Include their transverse perturbations needed
to compute all intersection numbers in a
compact set $K_i$ of that complement. These
are homology-cycle representatives; they
need not be embedded surfaces.

There is an actual orientation-preserving
smooth embedding of a neighbourhood of
$K_i$ into a nearby fibre of (C.PG.2).
Here is the local construction. On a
neighbourhood of $K_i$, $dg$ is surjective.
Choose smooth horizontal lifts of the two
coordinate fields of the real base disc;
a partition of unity preserves their
specified images under $dg$. Multiply by
a cutoff outside a slightly larger compact
neighbourhood and flow successively for
the real and imaginary parts of $s$.
For sufficiently small $s$, the trajectories
stay where the specified horizontal equations
hold. Their composition maps a neighbourhood
of $K_i$ in $E_i$ diffeomorphically to an
open subset of $g^{-1}(s)$. The image is
in that exact fibre by the two equations
for $dg$, and its orientation agrees with
the complex orientation by continuity
from $s=0$.

The finitely many compact sets $K_i$ lie
in different central components away
from all intersections. Choose their
ambient neighbourhoods disjoint, and
then choose the same sufficiently small
$s\ne0$ for all the flows. Their images
in the single torus $F=g^{-1}(s)$ remain
disjoint. Transverse intersection numbers
of the chosen cycles are preserved by
the local embeddings, and cycles from
different components have zero intersection.
Consequently their images give mutually
orthogonal positive subspaces of
$H_2(F,\mathbb R)$, of dimensions
$2p_g(E_i)$. Injectivity on each positive
subspace, and on their direct sum,
follows from preservation of its positive
intersection form.

The intersection form of an oriented
real four-torus is $3H$, where $H$ is
the hyperbolic plane. This follows
directly from the exterior-algebra
identification of its integral
cohomology: the three complementary
degree-two pairs give three hyperbolic
blocks. Its positive index is therefore
three. We have proved
\[
                    2\sum_i p_g(E_i)\le3.
                                                               \tag{C.PG.9}
\]
The sum is integral, so it is at most
one. Together with (C.PG.5) this is
(C.PG.1). The compact-core argument
used only multiplicity one, and hence
also proves the stated assertion for
any collection of such components in
a smooth finite-base-change model.
\end{proof}


% C.2.UN.73 from part_C2_compact_toric_gluing.tex

